\chapter{Conclusion}
\label{ch:conclusion}

Within the scope of this bachelor's thesis, a tiered classification and confidence-based safe-routing system was designed, implemented, and evaluated for pneumothorax detection in chest radiographs. The system pairs a lightweight MobileNetV2 screening stage with a heavier EfficientNet-B4 or Ark+ Swin stage, and routes each image to the appropriate tier according to an adaptively thresholded confidence estimate. On top of this cascade it layers three uncertainty mechanisms --- Monte Carlo Dropout, test-time augmentation, and conformal prediction --- together with temperature-scaling calibration, Stable-Diffusion synthetic augmentation of the minority class, and HiResCAM explanations.

The proposed architecture provides a practical answer to two fundamental problems in the clinical use of deep-learning models: their uniform and often excessive computational cost, and their inability to report calibrated uncertainty or to abstain. The fifteen-configuration ablation study showed that the lightweight tier resolves roughly three quarters of all cases on its own, that the uncertainty stack reduces the Expected Calibration Error from $0.074$ to $0.050$ without harming discrimination, and that the full system with an Ark+ Swin backbone and mixup/CutMix regularisation attains the best overall balance, with an AUC of $0.924$, the highest $F_1$ and Matthews correlation coefficient in the study, and a conformal predictor that meets its $95\%$ coverage target. A preliminary zero-shot evaluation on CheXpert indicated only a modest domain-shift penalty.

Equally important to the scientific result is the engineering result: the method is delivered as a complete, reproducible system, with a layered codebase, a FastAPI inference service, a React front end, Docker packaging, experiment tracking, and a continuous-integration pipeline that keeps every reported number honest and regenerable from the evaluation artifacts.

The work also makes its own limitations explicit --- the conservativeness of the conformal sets, the threshold sensitivity of the router, the restriction to a single binary pathology, and the preliminary scale of the cross-dataset evaluation --- and points to concrete next steps: a learned routing policy, adaptive conformal sets, a multi-pathology cascade, and a full-scale external evaluation. Taken together, this conformal-prediction-backed, uncertainty-aware tiered classification approach offers a strong and reproducible reference point for the design of efficient and trustworthy clinical decision-support systems.
